\article{重数一定理}{Multiplicity one theorems}{I. Piatetski-Shapiro}{1}{209--212}
\origsection{1}{两个定理}
设 $G=\GL_n$ 定义在整体域（global field）$k$ 上。本文讨论以下两个结果。
\begin{originaltheorem}{（1）重数一定理（Multiplicity one theorem）}
设 $\pi$ 是 $G(\A)$ 的不可约光滑容许表示（irreducible smooth admissible representation）。则 $\pi$ 在尖点形式（cusp form）空间中出现的重数至多为一。
\end{originaltheorem}
（2）任意不可约容许光滑表示 $\pi$ 都可写成 $\pi=\bigotimes_p\pi_p$，其中每个 $\pi_p$ 是局部群 $G_p$ 的不可约容许光滑表示。
\begin{originaltheorem}{强重数一定理（Strong multiplicity one theorem）}
设 $\pi_1=\bigotimes_p\pi_{1,p}$ 与 $\pi_2=\bigotimes_p\pi_{2,p}$ 是两个不可约表示；\footnote{\textbf{译者注（本次校订）：}这里必须假设二者都是尖点自守表示。对任意不可约容许表示，仅改变一个素位处的局部分量就能得到反例；下文证明也明确使用了尖点形式的 Fourier 展开。}假定对每个 $p\notin S$ 有 $\pi_{1,p}\cong\pi_{2,p}$，其中 $S$ 是有限集，当 $n>2$ 时还假定 $S$ 只含有限素位（finite place）。那么对所有 $p$ 都有 $\pi_{1,p}\cong\pi_{2,p}$。（从而 $\pi_1=\pi_2$。）
\end{originaltheorem}
先概述定理（1）的证明，所用的基本工具是 Whittaker 模型（Whittaker model）。先设 $k$ 为局部域（local field），$(\pi,V)$ 为不可约光滑表示，介绍其 Whittaker 模型。当 $k$ 为阿基米德域（archimedean）时，“光滑表示”指 $G$ 在某个 Hilbert 空间 $H$ 的 $C^\infty$ 向量空间 $V$ 上的表示，而 $G$ 在 $H$ 上酉作用（unitary action）；当 $k$ 非阿基米德（nonarchimedean）时，这个概念已在 Cartier 的讲座中介绍。令 $\psi$ 为 $k$ 的一个加法特征（additive character），\footnote{\textbf{译者注（本次校订）：}这里的 $\psi$ 须为非平凡加法特征，才能得到非退化的 Whittaker 模型及其唯一性；在后面的整体积分中，$\psi$ 则取为 $k\backslash\A$ 的非平凡特征。}令
\[
X=\left\{\begin{pmatrix}1&&*\\&\ddots&\\0&&1\end{pmatrix}\right\}
\]
为 $G$ 的标准极大幺幂子群（standard maximal unipotent subgroup）。那么 $(\pi,V)$ 的 Whittaker 模型 $W(\pi,\psi)$ 是某个非零交织映射 $V\to\Ind_X^G(\psi)$ 的像，其中当
\[
x=\begin{pmatrix}1&x_1&&*\\&1&\ddots&\\&&\ddots&x_{n-1}\\0&&&1\end{pmatrix}
\quad\text{时，}\quad \psi(x)=\psi(x_1+\cdots+x_{n-1}).
\]
具体地说，它由一族光滑函数 $\{W_v:G\to\C\mid v\in V\}$ 给出，满足：
\begin{enumerate}
\item[(i)] 对所有 $x\in X,g\in G$，$W_v(xg)=\psi(x)W_v(g)$；
\item[(ii)] 对所有 $g,h\in G$，$W_{\pi(h)v}(g)=W_v(gh)$。
\end{enumerate}
下面的唯一性定理是关键：非阿基米德情形由 Gelfand--Kazhdan 得到，一般局部域情形由 Shalika 得到（[1]、[2]）。
\begin{originaltheorem}{唯一性定理（Uniqueness theorem）}
对每个不可约容许光滑表示 $(\pi,V)$，至多存在一个 $W(\pi,\psi)$（其中 $\psi$ 固定）。
\end{originaltheorem}
当 $k$ 为阿基米德域时，假定 $(\pi,V)$ 可酉化（unitarizable），并且
\[
V=\{x\in H\mid (\mathcal Dx,\mathcal Dx)<\infty\text{，对包络代数（enveloping algebra）中的所有 }\mathcal D\}.
\]
这里 $H$ 是 $V$ 关于给定内积的完备化。还假定 $W_v(1)$ 关于半范数 $(\mathcal Dx,\mathcal Dx)$（$\mathcal D$ 属于包络代数）定义的拓扑，是 $V$ 上的连续线性泛函（continuous linear functional）。

整体 Whittaker 模型可按同样的方式定义，其唯一性也随即由上述局部结果推出。
\origsection{2}{整体 Fourier 分析（Global Fourier analysis）}
如前，设 $(\pi,V)$ 为容许不可约尖点表示，$\varphi\in V$。可定义
\[
W_\varphi(g)=\int_{X_k\backslash X_\A}\varphi(xg)\psi^{-1}(x)\,dx.
\]
整体 Fourier 分析表明，尖点形式由这个“Fourier 变换（Fourier transform）”唯一确定。在经典背景中，这是 Hecke 的结果；$n=2$ 时见 Jacquet--Langlands [3]；$n>2$ 时由 Piatetski-Shapiro [4] 与 Shalika [2] 独立得到。这一证明受到 S. I. Gelfand 在有限域上相应结果的启发 [5]。

由这些结果即可推出定理（1），因为
\[
\dim\Hom_G(V,W(\pi,\psi))=1\geq\dim\Hom_G(V,L_0^2).
\]
接着证明强重数一定理。先讨论 $n=2$ 的情形；需要下述引理。
\begin{originaltheorem}{小引理（Small lemma）}
（1）设 $k$ 为局部域，$(\pi_1,V_1),(\pi_2,V_2)$ 是具有 Whittaker 模型的两个不可约容许表示。那么存在 $v_1\in V_1,v_2\in V_2$，使
\[
W_{v_1}\begin{pmatrix}x&0\\0&1\end{pmatrix}
 =W_{v_2}\begin{pmatrix}x&0\\0&1\end{pmatrix},
\qquad W_{v_i}\in W(\pi_i,\psi).
\]
（2）若 $k=\R$ 或 $\C$，假定 $(\pi_1,H_1),(\pi_2,H_2)$ 为不可约无限维酉表示。以 $V_1$（$V_2$）表示 $H_1$（$H_2$）中所有光滑向量的集合。那么存在 $v_1\in V_1,v_2\in V_2$，使
\[
W_{v_1}\begin{pmatrix}x&0\\0&1\end{pmatrix}
 =W_{v_2}\begin{pmatrix}x&0\\0&1\end{pmatrix},
\qquad W_{v_i}\in W(\pi_i,\psi).
\]
\end{originaltheorem}
\begin{proof}
若 $k$ 为非阿基米德局部域，已知在 Kirillov 模型中，$V$ 包含 $k^*$ 上所有紧支集的 Schwartz--Bruhat 函数（Schwartz--Bruhat function），结论由此得到。\footnote{\textbf{译者注（本次校订）：}本引理及下文高维的对应引理需要所取向量、以及相等的限制函数均非零。取同一个非零紧支集函数即可保证这一点；若允许零函数，断言本身没有内容，也无法用于后面的证明。}

现在设 $k=\R$ 或 $\C$。Kirillov 定理（见 [8，第221页]）指出，$\GL(2,k)$ 的每个不可约无限维酉表示限制到子群
\[
P=\left\{\begin{pmatrix}a&b\\0&1\end{pmatrix}\right\}
\]
后仍不可约，因此作为 $P$ 的表示同构于 $P$ 的标准表示。故若 $\varphi(x)$ 是紧支集的 $C^\infty$ 函数，便存在 $v_1\in V_1,v_2\in V_2$，使
\[
W_{v_i}\begin{pmatrix}x&0\\0&1\end{pmatrix}=\varphi(x).
\]
\end{proof}
\paragraph*{注（Remark）}
若能将具有 Whittaker 模型的酉表示的内积写成与 $n=2$ 时类似的积分，就可以仿照 $n=2$ 的情形，对任意 $n$ 证明“小引理”，由此推出任意 $n$ 的强重数一定理。

接着给出 Jacquet--Langlands 对 $\GL(2,\A)$ 建立的从 Whittaker 模型恢复 $\varphi$ 的公式：
\[
\varphi(g)=\sum_{\lambda\in k^*}W_\varphi\left(\begin{pmatrix}\lambda&0\\0&1\end{pmatrix}g\right).\tag{*}
\]
现在设 $\pi_1,\pi_2$ 满足定理的假设。为证明结论，只需构造一个非零的 $\varphi\in V_1\cap V_2$，此时由 $(\pi_i,V_i)$ 的不可约性便有 $V_1=V_2$。而且由于 $B_k\backslash B_\A$ 在 $G_k\backslash G_\A$ 中稠密，只需构造两个在 $B_\A$ 上相等的非零函数 $\varphi_i\in V_i$（这里 $B$ 仍表示上三角矩阵群）。

由 Whittaker 模型的性质和 (*)，只需构造 Whittaker 函数 $W_1,W_2$，使对 $x\in\A^*$ 有
\[
W_1\begin{pmatrix}x&0\\0&1\end{pmatrix}=W_2\begin{pmatrix}x&0\\0&1\end{pmatrix}.
\]
可假设这些 $W_i$ 形如 $\prod_p W_i^p$；根据假设，只需在有限多个素位构造适当的 $W_i$，这些素位可逐一应用小引理。这类论证由 Shalika 在此地以及莫斯科方面独立发现。

对 $n\geq3$，需要一个类似的小引理（Gelfand--Kazhdan）：设 $k$ 为非阿基米德局部域，$(\pi_i,V_i)$（$i=1,2$）为具有 Whittaker 模型的不可约容许表示。存在 $v_i\in V_i$，使
\[
W_{v_1}\begin{pmatrix}h&0\\0&I\end{pmatrix}
 =W_{v_2}\begin{pmatrix}h&0\\0&I\end{pmatrix},
\qquad\text{对所有 }h\in\GL(n-1).
\]
再仿照 $n=2$ 的论证，就能用归纳法证明一般情形。作者预期，阿基米德局部域上也能证明这一引理；这样便可以去掉此前对 $S$ 不含无限素位（infinite place）的限制。

现在假设 $G$ 拟分裂（quasi-split），并满足如下传递性条件（transitivity condition）：$T(\A)$ 在
\[
\prod_{\alpha\text{ 为单根}}X_\alpha^*(\A)
\]
上传递作用。这里 $T$ 是一个 Borel 群中的极大 $k$ 环面（maximal $k$-torus），$X_\alpha^*=X_\alpha-\{I\}$，其中 $X_\alpha$ 是与单根（simple root）$\alpha$ 相应的根群（root group）。

若对所有 $\varphi\in V$ 都有
\[
W_\varphi(g)=\int_{X_k\backslash X_\A}\varphi(xg)\psi^{-1}(x)\,dx=0,
\]
则称自守尖点不可约表示 $(\pi,V)$ 为强尖点的（hypercuspidal），或退化尖点的（degenerate cuspidal）。从不含上半平面 $H=\{\operatorname{Im}z>0\}$ 因子的对称空间提升而来的全纯尖点形式，就属于这一类型。

若一个尖点自守形式与所有强尖点自守形式正交，就称它为一般型的（generic）；这里使用通常的标量积 $\int_{CG_k\backslash G_\A}\varphi\overline\psi\,dg$。

Howe 和作者在本次会议上给出的 Ramanujan 猜想的反例，是强尖点形式 [6]。作者仍认为 Ramanujan 猜想应保留其核心内容，并猜想它对一般型的不可约尖点自守表示成立。

现在概述一般型尖点自守形式的重数一定理的证明。局部 Whittaker 模型的唯一性定理仍然成立 [2]。虽然 Whittaker 函数不能唯一确定任意尖点形式，但由一般型的定义，它能唯一确定一般型尖点形式。因此，仿照前面的论证即可得到重数一。可以证明，每个拟分裂约化群上都存在一般型尖点形式。也可以证明，对这种群存在幺幂子群 $U$，使 $\int_{U_k\backslash U_\A}\varphi(ug)\,du$ 能够用 Whittaker 函数表示。由于除 $\GL(n)$ 外的任意群上都存在强尖点形式，单凭 Whittaker 函数便无法恢复任意 $\varphi$。有关细节将在 Novodvorskii 与作者即将发表的文章中给出。（笔记由 B. Seifert 和 L. Morris 整理。）

\begin{thebibliography}{8}
\bibitem[1]{c14-1} I. M. Gelfand and D. I. Kazhdan, \emph{Representation of the group $\GL(n,K)$, where $K$ is a local field}, Lie Groups and their Representations (Proc. Summer School of the Bolya-Janos Math. Soc., Budapest, 1971), Halsted, New York, 1975.
\bibitem[2]{c14-2} J. A. Shalika, \emph{The multiplicity one theorem for $\GL(n)$}, Ann. of Math. (2) 100 (1974).
\bibitem[3]{c14-3} H. Jacquet and R. P. Langlands, \emph{Automorphic forms on $\GL(2)$}, Lecture Notes in Math., vol. 114, Springer-Verlag, New York, 1970.
\bibitem[4]{c14-4} I. I. Piatetski-Shapiro, \emph{Euler subgroups}, Lie Groups and their Representations (Proc. Summer School of the Bolya-Janos Math. Soc., Budapest, 1971), Halsted, New York, 1975.
\bibitem[5]{c14-5} S. I. Gelfand, \emph{Representations of the general linear group over a finite field}, (Proc. Summer School of the Bolya-Janos Math. Soc., Budapest, 1971), Halsted, New York, 1975.
\bibitem[6]{c14-6} R. Howe and I. Piatetski-Shapiro, \emph{A counterexample to the ``generalized Ramanujan conjecture'' for (quasi-) split groups}, 本论文集，第1册，315--322页。
\bibitem[7]{c14-7} R. P. Langlands, \emph{On the notion of an automorphic representation}, 本论文集，第1册，203--207页。
\bibitem[8]{c14-8} I. M. Gelfand, M. I. Graev and I. Piatetski-Shapiro, \emph{Representation theory and automorphic functions}, W. B. Saunders Company, 1969.
\end{thebibliography}
\noindent 特拉维夫大学（University of Tel Aviv）。
\par\noindent\footnotesize 原文 AMS（MOS）学科分类（1970）：主分类 22E55；次分类 10D20。作者撰写本文期间受到 NSF 基金资助。版权 © 1979 American Mathematical Society。\normalsize
